\chapter{About the Case Study and Data Collection}

\section{Selection of Case Studies}
Six joints were selected: the knuckle joint, the socket and spigot cotter joint, the sleeve and cotter joint, the gib and cotter joint, the hex bolted joint and the protected type rigid flange coupling. Each has a published design procedure, and together they exercise tension, double shear, crushing, bending and torsion. The flange coupling is sized from a transmitted torque rather than an axial load, which confirms that the system is not limited to axially loaded joints. Their components, checks and duties are listed with the results in Table~\ref{tab:joints}.

\section{Data Collected}
For each joint three categories of data were collected: the empirical proportions as multiples of the governing dimension, every failure mode with its governing expression, and the worked example of the source. The allowable stresses and the preferred size series are given in Table~\ref{tab:allowables}; every computed diameter is rounded up to the next size in that series, as the sources direct.

\begin{table}[htbp]
\centering
\caption{Allowable stresses adopted, and the sizes to which computed diameters are rounded up}
\label{tab:allowables}
\small
\begin{tabularx}{\textwidth}{>{\raggedright\arraybackslash}p{3.6cm}>{\raggedright\arraybackslash}X>{\raggedright\arraybackslash}p{4.2cm}}
\toprule
Quantity & Value adopted & Source \\
\midrule
Allowable tensile stress & 75 MPa & Working value of the source \\
Allowable shear stress & 60 MPa & Working value of the source \\
Allowable crushing stress & 150 MPa & Working value of the source \\
Design factor & 1.25 & Sized for 1.25 times the load, reported at the load \\
Preferred diameters & 6 to 22 mm in steps of 2; 25, 30, 32, 35; 40 to 60 in steps of 5; then 70, 80, 90, 100 and larger & Standard shaft sizes, NPTEL Module~4 Lesson~2 \\
\bottomrule
\end{tabularx}
\end{table}

\section{Basis for Verification}
The worked examples of Bhandari, Example~4.2 for the cotter joint and Example~4.5 for the knuckle joint, and the NPTEL lesson each give the load, the allowable stresses and every resulting dimension. They were reproduced by the system and compared value by value; the outcome, including where the system and a source disagree, is reported in Chapter~7.
